% Options for packages loaded elsewhere
\PassOptionsToPackage{unicode}{hyperref}
\PassOptionsToPackage{hyphens}{url}
\documentclass[
  polish,
  11pt,
  a4paper,
]{article}
\usepackage{xcolor}
\usepackage[margin=2cm]{geometry}
\usepackage{amsmath,amssymb}
\setcounter{secnumdepth}{-\maxdimen} % remove section numbering
\usepackage{iftex}
\ifPDFTeX
  \usepackage[T1]{fontenc}
  \usepackage[utf8]{inputenc}
  \usepackage{textcomp} % provide euro and other symbols
\else % if luatex or xetex
  \usepackage{unicode-math} % this also loads fontspec
  \defaultfontfeatures{Scale=MatchLowercase}
  \defaultfontfeatures[\rmfamily]{Ligatures=TeX,Scale=1}
\fi
\usepackage{lmodern}
\ifPDFTeX\else
  % xetex/luatex font selection
\fi
% Use upquote if available, for straight quotes in verbatim environments
\IfFileExists{upquote.sty}{\usepackage{upquote}}{}
\IfFileExists{microtype.sty}{% use microtype if available
  \usepackage[]{microtype}
  \UseMicrotypeSet[protrusion]{basicmath} % disable protrusion for tt fonts
}{}
\makeatletter
\@ifundefined{KOMAClassName}{% if non-KOMA class
  \IfFileExists{parskip.sty}{%
    \usepackage{parskip}
  }{% else
    \setlength{\parindent}{0pt}
    \setlength{\parskip}{6pt plus 2pt minus 1pt}}
}{% if KOMA class
  \KOMAoptions{parskip=half}}
\makeatother
\usepackage{color}
\usepackage{fancyvrb}
\newcommand{\VerbBar}{|}
\newcommand{\VERB}{\Verb[commandchars=\\\{\}]}
\DefineVerbatimEnvironment{Highlighting}{Verbatim}{commandchars=\\\{\}}
% Add ',fontsize=\small' for more characters per line
\usepackage{framed}
\definecolor{shadecolor}{RGB}{248,248,248}
\newenvironment{Shaded}{\begin{snugshade}}{\end{snugshade}}
\newcommand{\AlertTok}[1]{\textcolor[rgb]{0.94,0.16,0.16}{#1}}
\newcommand{\AnnotationTok}[1]{\textcolor[rgb]{0.56,0.35,0.01}{\textbf{\textit{#1}}}}
\newcommand{\AttributeTok}[1]{\textcolor[rgb]{0.13,0.29,0.53}{#1}}
\newcommand{\BaseNTok}[1]{\textcolor[rgb]{0.00,0.00,0.81}{#1}}
\newcommand{\BuiltInTok}[1]{#1}
\newcommand{\CharTok}[1]{\textcolor[rgb]{0.31,0.60,0.02}{#1}}
\newcommand{\CommentTok}[1]{\textcolor[rgb]{0.56,0.35,0.01}{\textit{#1}}}
\newcommand{\CommentVarTok}[1]{\textcolor[rgb]{0.56,0.35,0.01}{\textbf{\textit{#1}}}}
\newcommand{\ConstantTok}[1]{\textcolor[rgb]{0.56,0.35,0.01}{#1}}
\newcommand{\ControlFlowTok}[1]{\textcolor[rgb]{0.13,0.29,0.53}{\textbf{#1}}}
\newcommand{\DataTypeTok}[1]{\textcolor[rgb]{0.13,0.29,0.53}{#1}}
\newcommand{\DecValTok}[1]{\textcolor[rgb]{0.00,0.00,0.81}{#1}}
\newcommand{\DocumentationTok}[1]{\textcolor[rgb]{0.56,0.35,0.01}{\textbf{\textit{#1}}}}
\newcommand{\ErrorTok}[1]{\textcolor[rgb]{0.64,0.00,0.00}{\textbf{#1}}}
\newcommand{\ExtensionTok}[1]{#1}
\newcommand{\FloatTok}[1]{\textcolor[rgb]{0.00,0.00,0.81}{#1}}
\newcommand{\FunctionTok}[1]{\textcolor[rgb]{0.13,0.29,0.53}{\textbf{#1}}}
\newcommand{\ImportTok}[1]{#1}
\newcommand{\InformationTok}[1]{\textcolor[rgb]{0.56,0.35,0.01}{\textbf{\textit{#1}}}}
\newcommand{\KeywordTok}[1]{\textcolor[rgb]{0.13,0.29,0.53}{\textbf{#1}}}
\newcommand{\NormalTok}[1]{#1}
\newcommand{\OperatorTok}[1]{\textcolor[rgb]{0.81,0.36,0.00}{\textbf{#1}}}
\newcommand{\OtherTok}[1]{\textcolor[rgb]{0.56,0.35,0.01}{#1}}
\newcommand{\PreprocessorTok}[1]{\textcolor[rgb]{0.56,0.35,0.01}{\textit{#1}}}
\newcommand{\RegionMarkerTok}[1]{#1}
\newcommand{\SpecialCharTok}[1]{\textcolor[rgb]{0.81,0.36,0.00}{\textbf{#1}}}
\newcommand{\SpecialStringTok}[1]{\textcolor[rgb]{0.31,0.60,0.02}{#1}}
\newcommand{\StringTok}[1]{\textcolor[rgb]{0.31,0.60,0.02}{#1}}
\newcommand{\VariableTok}[1]{\textcolor[rgb]{0.00,0.00,0.00}{#1}}
\newcommand{\VerbatimStringTok}[1]{\textcolor[rgb]{0.31,0.60,0.02}{#1}}
\newcommand{\WarningTok}[1]{\textcolor[rgb]{0.56,0.35,0.01}{\textbf{\textit{#1}}}}
\usepackage{longtable,booktabs,array}
\usepackage{caption}
\captionsetup[table]{skip=6pt}
\newcounter{none} % for unnumbered tables
\usepackage{calc} % for calculating minipage widths
% Correct order of tables after \paragraph or \subparagraph
\usepackage{etoolbox}
\makeatletter
\patchcmd\longtable{\par}{\if@noskipsec\mbox{}\fi\par}{}{}
\makeatother
% Allow footnotes in longtable head/foo
\IfFileExists{footnotehyper.sty}{\usepackage{footnotehyper}}{\usepackage{footnote}}
\makesavenoteenv{longtable}
\usepackage{graphicx}
\makeatletter
\newsavebox\pandoc@box
\newcommand*\pandocbounded[1]{% scales image to fit in text height/width
  \sbox\pandoc@box{#1}%
  \Gscale@div\@tempa{\textheight}{\dimexpr\ht\pandoc@box+\dp\pandoc@box\relax}%
  \Gscale@div\@tempb{\linewidth}{\wd\pandoc@box}%
  \ifdim\@tempb\p@<\@tempa\p@\let\@tempa\@tempb\fi% select the smaller of both
  \ifdim\@tempa\p@<\p@\scalebox{\@tempa}{\usebox\pandoc@box}%
  \else\usebox{\pandoc@box}%
  \fi%
}
% Set default figure placement to htbp
\def\fps@figure{htbp}
\makeatother
\ifLuaTeX
\usepackage[bidi=basic,shorthands=off]{babel}
\else
\usepackage[bidi=default,shorthands=off]{babel}
\fi
\ifLuaTeX
  \usepackage{selnolig} % disable illegal ligatures
\fi
\setlength{\emergencystretch}{3em} % prevent overfull lines
\providecommand{\tightlist}{%
  \setlength{\itemsep}{0pt}\setlength{\parskip}{0pt}}
\usepackage{fontspec}
\defaultfontfeatures{Ligatures=TeX}
\usepackage{amsmath,amssymb}
\usepackage{booktabs}
\usepackage{array}
\usepackage{geometry}
\usepackage{enumitem}
\usepackage{xcolor}
\usepackage{tcolorbox}
\usepackage{fancyhdr}
\usepackage{titlesec}
\usepackage{multicol}

\geometry{margin=1.8cm, top=2.5cm, bottom=2cm}
\pagestyle{fancy}
\fancyhf{}
\rhead{\small Projektowanie badań — 2026/27}
\lhead{\small Zarządzanie informacją | ISI UJ}
\cfoot{\thepage}

\titleformat{\section}{\large\bfseries\color{blue!70!black}}{\thesection.}{0.5em}{}
\titleformat{\subsection}{\normalsize\bfseries}{\thesubsection.}{0.5em}{}
\titleformat{\subsubsection}{\small\bfseries\color{gray!70!black}}{\thesubsubsection.}{0.5em}{}
\titlespacing*{\section}{0pt}{1.5ex}{1ex}
\titlespacing*{\subsection}{0pt}{1ex}{0.5ex}

\setlist[itemize]{nosep, leftmargin=1.5em}
\setlist[enumerate]{nosep, leftmargin=1.5em}

\tcbset{boxrule=0.5pt, left=4pt, right=4pt, top=4pt, bottom=4pt, fonttitle=\bfseries\small}

\newtcolorbox{definicja}[1][]{colback=blue!5, colframe=blue!50!black, title={#1}}
\newtcolorbox{uwagabox}[1][]{colback=orange!10, colframe=orange!70!black, title={#1}}
\newtcolorbox{wzorbox}[1][]{colback=gray!5, colframe=gray!50!black, title={#1}}
\newtcolorbox{przykladbox}[1][]{colback=purple!5, colframe=purple!50!black, title={#1}}

\usepackage{bookmark}
\IfFileExists{xurl.sty}{\usepackage{xurl}}{} % add URL line breaks if available
\urlstyle{same}
% fallback for those not using the hyperref driver hyperxmp:
\makeatletter
\@ifundefined{xmpquote}{\newcommand{\xmpquote}[1]{#1}}{}
\makeatother
\hypersetup{
  pdftitle={Handout W01 --- Pytania{,} projekty badań i skale pomiarowe},
  pdfauthor={Marek Deja},
  pdflang={pl-PL},
  hidelinks,
  pdfcreator={LaTeX via pandoc}}

\title{Handout W01 --- Pytania, projekty badań i skale pomiarowe}
\usepackage{etoolbox}
\makeatletter
\providecommand{\subtitle}[1]{% add subtitle to \maketitle
  \apptocmd{\@title}{\par {\large #1 \par}}{}{}
}
\makeatother
\subtitle{Zarządzanie informacją --- część ilościowa --- 2026/27}
\author{Marek Deja}
\date{2026/27}

\begin{document}
\maketitle

\section{W01 --- pytanie i zakres
dowodu}\label{w01-pytanie-i-zakres-dowodu}

Problem instytucji → pytanie → jednostka/populacja/próba → pomiar →
analiza → ograniczony wniosek. Średnia w próbie jest opisem; wniosek o
populacji wymaga projektu i założeń.

\section{Schematy badań}\label{schematy-badaux144}

{\def\LTcaptype{none} % do not increment counter
\begin{longtable}[]{@{}
  >{\raggedright\arraybackslash}p{(\linewidth - 2\tabcolsep) * \real{0.5000}}
  >{\raggedright\arraybackslash}p{(\linewidth - 2\tabcolsep) * \real{0.5000}}@{}}
\toprule\noalign{}
\begin{minipage}[b]{\linewidth}\raggedrigh
Schema
\end{minipage} & \begin{minipage}[b]{\linewidth}\raggedrigh
Kluczowe pytanie
\end{minipage}
\midrule\noalign{}
\endhead
\bottomrule\noalign{}
\endlastfoo
Obserwacja przekrojowa & Czy cechy współwystępują w danym okresie?
Obserwacja podłużna & Czy wynik zmienia się u tych samych osób?
Interwencja przed/po & Czy zmiana mogła mieć inną przyczynę?
Randomizacja & Czy przydzielono interwencję losowo?
Quasi-eksperyment/naturalny & Co uzasadnia porównywalność i
ekspozycję?
Synteza & Czy badania są odpowiednie, porównywalne i dobrej jakości?
\end{longtable}
}

Piramida dowodu zależy od pytania. Randomizacja przydziału i losowanie
próby są różnymi operacjami. Duże N nie usuwa selekcji ani błędnego
pomiaru.

\section{Pytanie i hipoteza}\label{pytanie-i-hipoteza}

„Jak rozkłada się czas znalezienia dokumentu wśród aktywnych
użytkowników w semestrze zimowym?'' --- pytanie opisowe. „Czy indeks
dostępności różni się między nowymi i doświadczonymi?'' --- porównanie.
\(H_0:\mu_1-\mu_2=0\), \(H_1:\mu_1-\mu_2\ne0\). Kierunek alternatywy
wybiera się przed wynikiem.

\section{Skale i jednostki}\label{skale-i-jednostki}

{\def\LTcaptype{none} % do not increment counter
\begin{longtable}[]{@{}
  >{\raggedright\arraybackslash}p{(\linewidth - 4\tabcolsep) * \real{0.3333}}
  >{\raggedright\arraybackslash}p{(\linewidth - 4\tabcolsep) * \real{0.3333}}
  >{\raggedright\arraybackslash}p{(\linewidth - 4\tabcolsep) * \real{0.3333}}@{}}
\toprule\noalign{}
\begin{minipage}[b]{\linewidth}\raggedrigh
Skala
\end{minipage} & \begin{minipage}[b]{\linewidth}\raggedrigh
Przykład
\end{minipage} & \begin{minipage}[b]{\linewidth}\raggedrigh
Uprawnione relacje
\end{minipage}
\midrule\noalign{}
\endhead
\bottomrule\noalign{}
\endlastfoo
Nominalna & grupa/typ zasobu & równość, liczebności
Porządkowa & pojedyncza pozycja 1--5 & dodatkowo porządek
Przedziałowa & rok kalendarzowy & dodatkowo różnica
Ilorazowa & czas w minutach/liczba pobrań & dodatkowo stosunek i
naturalne zero
\end{longtable}
}

Klasa R nie zastępuje skali. Kod kategorii pozostaje etykietą. Indeks
wymaga treści, kierunku pozycji i reguły braków.

\section{Sprawdź na przykładzie}\label{sprawdux17a-na-przykux142adzie}

Pięć czasów opisuje małą grupę. Średnia i mediana odpowiadają na różne
pytania o centrum. Odczytaj N i jednostkę, a potem wskaż brakującą
informację o doborze.

\begin{Shaded}
\begin{Highlighting}[]
\NormalTok{x }\OtherTok{\textless{}{-}} \FunctionTok{c}\NormalTok{(}\DecValTok{6}\NormalTok{, }\DecValTok{9}\NormalTok{, }\DecValTok{4}\NormalTok{, }\DecValTok{12}\NormalTok{, }\DecValTok{8}\NormalTok{)}
\FunctionTok{c}\NormalTok{(}\AttributeTok{N =} \FunctionTok{length}\NormalTok{(x), }\AttributeTok{srednia =} \FunctionTok{mean}\NormalTok{(x), }\AttributeTok{mediana =} \FunctionTok{median}\NormalTok{(x))}
\end{Highlighting}
\end{Shaded}

\begin{verbatim}
##       N srednia mediana
##     5.0     7.8     8.0
\end{verbatim}

\section{Własny szkic i notatki}\label{wux142asny-szkic-i-notatki}

Problem:
\ldots\ldots\ldots\ldots\ldots\ldots\ldots\ldots\ldots\ldots\ldots\ldots{}
Populacja/okres:
\ldots\ldots\ldots\ldots\ldots\ldots\ldots\ldots\ldots\ldots\ldots{}

Jednostka:
\ldots\ldots\ldots\ldots\ldots\ldots\ldots\ldots\ldots\ldots\ldots.
Zmienna i skala:
\ldots\ldots\ldots\ldots\ldots\ldots\ldots\ldots\ldots\ldots\ldots.

Pytanie:
\ldots\ldots\ldots\ldots\ldots\ldots\ldots\ldots\ldots\ldots\ldots\ldots{}
Co mogłoby podważyć hipotezę?
\ldots\ldots\ldots\ldots\ldots\ldots\ldots{}

Jedno ograniczenie doboru lub pomiaru:
\ldots\ldots\ldots\ldots\ldots\ldots\ldots\ldots\ldots\ldots\ldots\ldots\ldots\ldots\ldots\ldots\ldots\ldots\ldots..

Ocena: zadania 50\% (każde 5\%), projekt 50\%; indywidualny plan i
analiza ilościowa, bez egzaminu.

\end{document}
